\section{Causal Effect Confirmation}

The third stage turns each frozen hypothesis into a family-valid randomized
experiment on held-out tasks. Typed variants and their controls are completed
before assignment, and analysis estimates the effect of assigned prompt roles
without selecting cases by realized edit success.

\subsection{Typed Intervention Protocols}

A \texttt{TargetSpec} binds the feature family and catalog identity, ADD or
REMOVE, source prompt role, and exact counterpart requirements. ADD requires
an applicable, resolved \textsc{absent} source state. REMOVE requires a
provenance-bound \textsc{present} clause and a pre-attested, task-preserving
neutral counterpart. Execution is text-native or graph-native and uses either
a deterministic renderer or a locked LLM executor; the run fixes this choice
and its complete provenance.

Four-arm protocols are specific to Safety ADD and Safety REMOVE:

\begin{itemize}
  \item Safety ADD uses \path{TARGET_PATCH}, \path{NOOP_REWRITE},
  \path{LENGTH_MATCHED_PLACEBO}, and
  \path{GENERIC_SECURITY_REMINDER}.
  \item Safety REMOVE uses \path{TARGET_REMOVE}, \path{NOOP_RETAIN},
  \path{LENGTH_MATCHED_SHAM_EDIT}, and
  \path{GENERIC_SECURITY_REPLACEMENT}.
\end{itemize}

ADD and REMOVE remain separate candidate slots, hypotheses, protocols,
contrasts, multiplicity coordinates, and evidence. REMOVE may only restore an
attested neutral counterpart and never inserts an unsafe request; ADD and
REMOVE therefore cannot be combined in one confirmable hypothesis.

Other families retain family-specific protocols. Task-function ADD/REMOVE
requires target, no-op, and length-matched placebo roles; a fourth active
control is permitted only with a reviewed catalog entry and independent
functional outcome. Presentation experiments use target and no-op plus an
optional matched presentation control and serve as negative controls.

Each arm has an explicit \texttt{AllowedDelta} that bounds its declared target
or control change and fixes protected task, safety, and presentation
projections. Exact source/counterpart provenance and the
\texttt{SecurityNeutralPromptInvariant} ensure that no variant claims a
vulnerability, requests unsafe behavior, or reveals an expected Oracle label
or effect direction.

\subsection{Variant Freeze and Block Randomization}

For each $h$, the protocol first freezes a finite task-independent realization
set and probabilities $Q_h$. For every eligible task and every realization,
one task-specific bundle contains all four arm texts and exact digests. Every
candidate arm is independently re-extracted with the run-locked Prompt
extractor while it is blind to the arm role and target. Before assignment, hard
gates validate schema and catalog integrity, exact evidence and provenance,
canonical round trips, security-neutral task preservation, context/task/non-
target invariance, the arm-specific \texttt{AllowedDelta}, and absence of
downstream information. A task enters the common-support population only when
every required arm and realization passes before outcomes exist; failed
realizations are never deleted or renormalized.

Target change, semantic compliance, and permissible non-target drift are
recorded during this arm-blind validation as intervention-realization
diagnostics. They
do not decide the hard-gate result, and \texttt{target\_changed}, semantic
compliance, and non-target drift never filter the primary assigned-arm ITT
denominator. If one arm fails a hard gate, the complete task--hypothesis bundle
is reported as a pre-randomization exclusion and receives no assignment.

Within each complete frozen block, a manifested deterministic RNG assigns
request-randomness slots to a balanced permutation of the protocol's finite arm
roles. The canonical block binds semantic cluster, task instance, hypothesis,
target, global realization specification, task realization bundle, model, and
arm protocol. The assignment manifest additionally binds prompt variant,
nullable provider seed, model parameters, RNG identity, request order, and
request ID before generation. Mutation of any canonical coordinate changes the
block identity, and provider execution order cannot alter treatment.

\subsection{Outcomes and Assigned-Arm ITT}

Generated code enters only the independent, versioned security Oracle and
functional evaluator. We decompose code production and Oracle support as
$Y_C$ and $Y_E$. The primary oracle-evaluable secure-code yield is
$Y_CY_E I(\text{Oracle secure})$; secure-and-functional joint success further
requires the frozen functional check to pass and is the key practical
secondary outcome. A valid terminal generation failure or syntactically
invalid result has $Y_C=0,Y_E=0$; valid but Oracle-unsupported or unknown code
has $Y_C=1,Y_E=0$. These remain observed assigned-arm outcomes under the
frozen encoding. Missing, corrupt, or unavailable required producer evidence
is not encoded as an outcome; it blocks publication and triggers deterministic
replay of the same manifest.

For frozen hypothesis $h$, target $q$, operation $o$, model $m$, and outcome
$Y$, the primary contrast is the hypothesis-specific target-minus-no-op
assigned-arm ITT risk difference

\[
  \Delta^{\mathrm{ITT}}_{h,q,o,m}
  = \mathbb{E}[Y \mid A=\mathrm{target}]
  - \mathbb{E}[Y \mid A=\mathrm{no\mbox{-}op}].
\]

Every committed randomized assignment enters the denominator under its
assigned role. Point estimates average over the frozen task population and
realization distribution $Q_h$. Uncertainty resamples semantic task clusters
and carries all descendant hypotheses, models, arms, realizations, and request
slots together. Simultaneous intervals follow the frozen hypothesis/model/
contrast family. Target-minus-placebo and, where family-valid,
target-minus-generic-control contrasts assess specificity. Effects remain
separate by hypothesis, target, operation, model, and CWE unless another
estimand was approved in advance. Task improvement, harm, and tie rates remain
descriptive diagnostics, not individual causal flips or substitutes for ITT.

\subsection{Evidence Levels}

\model uses the prospective evidence levels below:

\begin{enumerate}
  \item \textbf{Observational Candidate}: a discover-split stable possible path
  frozen before confirmation, without a causal confirmation claim.
  \item \textbf{Randomized Policy Effect}: the preregistered
  target-minus-no-op ITT follows the frozen expected direction, meets the
  independent-task requirement, and has a multiplicity-adjusted interval
  excluding zero with complete provenance.
  \item \textbf{Target-Specific Policy Effect}: a randomized effect also supported by
  the preregistered target-minus-placebo and, where family-valid,
  target-minus-generic contrasts.
  \item \textbf{Realization-Robust Policy Effect}: the average policy effect
  is confirmed and all frozen direction, equivalence, interaction, and
  leave-one-realization-out conditions hold simultaneously.
  \item \textbf{Cross-Model Replication}: separately estimated model effects
  meet the frozen replication rule without creating a pooled universal effect.
  \item \textbf{Bidirectional Support}: separate ADD and REMOVE protocols
  independently confirm opposite expected directions for the same feature.
  \item \textbf{Directionally Consistent but Inconclusive}: the point estimate
  follows the frozen direction but its adjusted interval includes zero or
  independent-task support is insufficient.
  \item \textbf{Null / Conflicting / Non-Evaluable}: the preregistered analysis
  respectively detects no effect, estimates the opposite direction, or cannot
  publish an effect because no valid randomization universe or required
  infrastructure/provenance exists.
\end{enumerate}

Valid post-assignment terminal non-successes remain assigned data and do not
create a non-evaluable result. Discovery, demo, and smoke-test artifacts cannot
be promoted to confirmatory evidence.
